% Methods text only; see PAPER_IMPLEMENTATION.md for claims and evidence limits.
% Bibliography: PAPER_REFERENCES.bib
We implement a Gemmini backend in a fork of Apache TVM v0.19.0, following TVM's
separation of graph and operator optimization \cite{chen2018tvm}.
The backend uses TensorIR and TVM's TensorIntrin interface: each intrinsic pairs
a mathematical computation description with a target implementation, and
schedule-driven tensorization replaces matching computation blocks
\cite{feng2023tensorir}.
The same pinned Gemmini software specification supplied to Merlin defines the
backend's capability and numerical requirements. Intrinsic recognition and
finite tests establish implementation coverage; they do not discharge every
specification proof obligation or qualify a deployed bitstream.

Graph preparation uses stock Relax legalization, constant folding, binding
canonicalization, operator/TIR fusion, reshape rewriting and dead-code
elimination. Common-subexpression elimination is an optional stock graph choice;
the initial configuration uses fusion level 2 with CSE disabled. Target schedule
rules compose ordinary TensorIR transformations for tiling, loop ordering,
blockization, cache placement and layout. Their traced choices expose intrinsic
tile dimensions, operand-reuse groups, dataflow, staging and legal
hoisting/pipelining alternatives. A capacity-checked, shape-derived schedule
provides the initial baseline. Necessary target-specific integration supplies
workload annotation, intrinsic implementations, final schedule placement, SRAM
allocation and DMA synchronization, together with a primitive command runtime
and static RV64 packaging. Explicit accumulator initialization may tensorize;
automatic bias relocation, constant-weight prepacking and scale/readout
normalization passes are excluded.

We use MetaSchedule's search-space and measurement interfaces
\cite{shao2022metaschedule},
with stock ReplayTrace and no learned cost model. The configured allowance is
at most 25 schedule proposals per kernel case or ResNet owner chunk, shared
across its internal tasks and graph alternatives. Completed correct warm-cycle
measurements select schedules; cycles are converted to seconds at the TVM
Runner boundary. Measurements do not guide ReplayTrace sampling. Three baseline
timings calibrate a median-based wall watchdog; failed or timed-out attempts do
not supply completed cycle labels. Matching traces are applied through
JSONDatabase and MetaScheduleApplyDatabase, retaining the measured default for
misses or unsuccessful search. Final comparisons require rebuilding and timing
the complete kernel or connected model, including executed host work, packing,
DMA and completion; sums of isolated operator minima are not E2E results.
